\section{共同编年：九六〇至一一二七年}

以下先按时间推进北宋已有的可用材料。辽、西夏与金在这条时间线上只在已经进入北宋事件链、且材料明确支持相互关系处出现；它们不是宋朝的附庸背景。编年之后的各节再把国家、边疆、财政、社会与文化问题交织起来阅读，避免用一个后来的结局替代当时各方的选择。

\section{从东京出发}

\begin{event}{\eventref{NS-0960-CHENQIAO}{陈桥建宋}（九六〇年正月；后周禁军与赵匡胤集团；陈桥、开封；建国可核，军中密议与动机未能复原）}
《宋史·太祖本纪》《续资治通鉴长编》建隆元年正月条及《宋会要辑稿·礼》可以交叉追踪即位、年序与礼仪语言。它们成书层次不同；“黄袍加身”的仪式细节和全体参与者的意图，不能据后出的铺陈写成现场。
\end{event}

公元960年正月，赵匡胤在陈桥兵变后即位，国号宋。\hyperlink{evt:NS-0960-CHENQIAO}{陈桥建宋（960）}新朝首先占有的，并不只是一顶皇帝的冠冕，而是开封这座已经运转的都城：宫城、官署、仓储、汴河和驻守其间的禁军，把政令、粮食与兵员系在同一处。东京城遗址所见的城垣、街道和水系，可以帮助理解这种承继；它并不容许我们把北宋都城的每一层空间都写得仿佛亲见。宋廷承用开封为东京，正因为这里既是后周的中枢，也是维持百官与军队供给的节点。

\phantomsection\label{fig:vol04:northern-song:founding-chenqiao-kaifeng}
\begin{center}
\begin{tikzpicture}[x=.88cm,y=.88cm,font=\small]
  \fill[paper] (0,0) rectangle (11.6,6.25);

  \fill[source!16] (.55,3.85) -- (3.65,4.15) -- (4.1,5.75) -- (.8,5.75) -- cycle;
  \node[align=center] at (2.25,5.0) {北方行军地带\\后周禁军北伐动员};
  \fill[timeline!23] (.75,.7) -- (4.35,.85) -- (4.1,3.15) -- (.55,3.0) -- cycle;
  \node[align=center] at (2.35,1.9) {陈桥\\军队驻泊与政权转换};
  \fill[dynasty!24] (7.1,1.0) -- (10.95,1.2) -- (11.05,5.35) -- (7.45,5.2) -- cycle;
  \node[align=center] at (9.1,3.65) {开封（东京）\\宫城、官署、仓储\\与禁军供给节点};

  \draw[very thick,color=source!75] (4.2,.6) .. controls (5.7,1.8) and (6.2,4.45) .. (7.1,5.85);
  \node[text=source,rotate=55,above,font=\scriptsize] at (5.55,3.25) {汴河—都城水运轴};
  \draw[thick,dashed,color=timeline!85] (4.0,4.55) .. controls (5.55,4.5) and (6.25,3.95) .. (7.25,3.85);
  \node[text=timeline!85!black,above,font=\scriptsize] at (5.7,4.3) {陈桥至开封方向};

  \node[draw=timeline!80,fill=paper,rounded corners=2pt,align=center,inner sep=3pt] (chenqiao)
    at (3.0,2.55) {960年正月\\陈桥政权转换};
  \node[draw=dynasty!80,fill=paper,rounded corners=2pt,align=center,inner sep=3pt] (kaifeng)
    at (7.8,2.25) {承用后周中枢\\改国号宋};
  \draw[->,ultra thick,color=dynasty] (chenqiao.east) -- node[above,align=center,font=\scriptsize]
    {入据京师\\取得既有官僚与仓储} (kaifeng.west);

  \node[draw=source!70,fill=paper,rounded corners=2pt,align=center,inner sep=3pt] (supply)
    at (9.2,5.0) {政令、粮饷、兵员\\向东京汇集};
  \draw[->,thick,color=source] (supply.south west) -- (8.2,4.35);
  \node[text=source,align=center,font=\scriptsize] at (5.6,.85)
    {建国不只是一场仪式；都城、河运与禁军供给也须被接续};
\end{tikzpicture}

\smallskip
\small\textit{图\quad 陈桥—开封与北宋建国中枢示意：据官修纪传《宋史》卷1〈太祖本纪一〉、南宋编年汇集《续资治通鉴长编》建隆元年正月条、后世辑存的制度汇编《宋会要辑稿·礼》《宋会要辑稿·方域》，并参照开封北宋东京城遗址的城垣、街道与水系考古资料，说明政权转换与承用京师、仓储和水运节点的关系。黄袍仪式细节及都城空间承继层次均须谨慎。图中政权边界、水道和陈桥—开封路线均为约略示意。}
\end{center}


建国的日期和权力转换的基本过程，可由《宋史》太祖本纪与《续资治通鉴长编》建隆元年条相互勾连；建国礼仪则还须参看《宋会要辑稿·礼》。三种材料回答的不是同一个问题。《宋史》是元末官修史，给出后来王朝整理出的帝纪框架；《长编》由南宋李焘编年汇集宋廷旧材料，保存了更密的政事线索，却也带着后来的取舍；《宋会要辑稿》是自《永乐大典》辑出的制度条目，适合追问礼仪和机构，不能把它当作现场实录。后世最熟悉的“黄袍加身”图景，以及围绕陈桥的许多仪式性细节，正是在这些不同层次的叙述中逐渐定型。能够确定的是政权已经转换；不能据较晚的铺陈，补写帐中密语或人人共有的意图。

\begin{sourcenote}{陈桥能证实到哪里}
《宋史》《续资治通鉴长编》与《宋会要辑稿》可分别校读建国顺序、政事线索与礼仪语言；它们并非同一时代、同一体裁的现场记录。政权转换可以成立，军中密议和所有参与者的动机则不能由后出的铺陈补足。
\end{sourcenote}

宋初最急迫的功课，是使取得京城不再等于又一次禁军夺位。961年前后，太祖以任官、调动和优礼宿将等办法，削弱禁军统帅各自掌握的独立兵权。后世把它收束为一次饮宴中的“杯酒释兵权”，但《宋史》太祖本纪、石守信传与《长编》建隆二年条所能支持的，是一个持续的权力重组，而非一段可以逐字复原的对话。将领的个人声望被纳入任命与调动，军队的供养和指挥愈来愈回到东京；文官监督、轮调与财赋供给共同构成约束。它降低了五代式统兵者凭一支禁军改朝换代的风险，也使国家安全更依赖首都能够持续征粮、转运和支付。

\section{把割据地带接入中枢}

这种能力随后被用于统一，却不应把统一理解为一张地图在数年间自然合拢。963年宋军入荆南，结束高氏政权，取得长江中游的要点；《宋史》太祖本纪、《长编》乾德元年条和《宋会要辑稿·方域》的相关条目，让我们看到征服、改置州县和区域接收是相连的过程。荆南处在宋、后蜀和南唐之间，体量有限而位置关键；它的归入为沿江行动打开了较可靠的节点，并不意味着地方的旧有官员、运输和赋役立即换成同一种秩序。

964年冬伐后蜀，次年孟昶降宋，四川遂入宋境。成都平原的资源使西南具有吸引力，剑门道路和山地运输却提醒进军者：胜利后仍有粮道、驻军、降附者和旧官安置的问题。《宋史》的太祖本纪与后蜀世家提供了宋方的事件骨架，《长编》乾德二、三年条补足年序；军纪、赏赐的范围和地方社会的感受，则不能只由征服者的叙述断定。统一在这里同时是一项行政接收工程。后来蜀地的赋役与控制压力会以新的形式浮现，不能倒过来把965年的每一项安排都预言为必然失败。

岭南和江南的进入，进一步显示财政、交通与战争相互借力。971年宋军攻取广州，南汉刘鋹降服；宋朝由此接触岭南的港口、盐业和既有海上贸易网络，须在路州体系之外处理边地军政。974年渡江围金陵、975年李煜降宋，则把长江下游的攻城、漕运和江南财赋集中到同一场战争中。《宋史》南汉、南唐世家和太祖本纪同《长编》开宝四年、七年至八年条，支持这些节点的次序；战报里的兵数、伤亡和俘获数字却不能未经核对便化作确定的规模。金陵陷落后，宋廷获得江南的税粮、手工业和水网地区，但它面对的是接收而非空地：如何让旧政权辖下的官民、仓储与税源进入新朝的账目，才是统一得以延续的条件。

\phantomsection\label{fig:vol04:northern-song:early-state-capacity-transport}
\begin{center}
\begin{tikzpicture}[x=.88cm,y=.88cm,font=\small]
  \fill[paper] (0,0) rectangle (11.8,6.35);

  \fill[source!16] (.55,3.7) -- (3.55,3.5) -- (3.75,5.85) -- (.8,5.95) -- cycle;
  \node[align=center] at (2.12,4.8) {河东、河北\\边防与军需压力};
  \fill[timeline!21] (.7,.65) -- (4.0,.75) -- (3.7,2.9) -- (.5,3.0) -- cycle;
  \node[align=center] at (2.2,1.78) {四川\\山地道路、粮赋\\与接收后的治理};
  \fill[dynasty!19] (8.3,.65) -- (11.25,.8) -- (11.15,3.2) -- (8.1,2.95) -- cycle;
  \node[align=center] at (9.7,1.8) {江南\\水网、税粮、手工业\\与渡江战争};
  \fill[timeline!18] (8.1,3.65) -- (11.25,3.55) -- (11.3,5.9) -- (8.35,5.8) -- cycle;
  \node[align=center] at (9.85,4.85) {岭南\\港口、盐业与\\边地军政};

  \node[draw=dynasty!85,fill=dynasty!15,rounded corners=3pt,align=center,inner sep=5pt] (tokyo)
    at (5.85,3.15) {东京开封\\官署、仓储、禁军\\与文书调度};
  \node[draw=source!70,fill=paper,rounded corners=2pt,align=center,inner sep=3pt] (officials)
    at (5.85,5.18) {任命、轮调、文官监督\\与地方接收};
  \draw[<->,thick,color=source] (tokyo.north) -- node[right,font=\scriptsize]{调度、考核} (officials.south);

  \draw[->,ultra thick,color=source!85] (3.6,4.55) -- node[above,sloped,font=\scriptsize]{兵员、粮饷、信息} (5.15,3.65);
  \draw[->,ultra thick,color=timeline!90] (3.55,2.05) -- node[below,sloped,font=\scriptsize]{道路、赋役与接收} (5.12,2.75);
  \draw[->,ultra thick,color=dynasty!88] (8.05,1.95) -- node[below,sloped,font=\scriptsize]{水路、税粮与军需} (6.55,2.76);
  \draw[->,ultra thick,color=timeline!85] (8.05,4.65) -- node[above,sloped,font=\scriptsize]{港口、盐业与政务} (6.55,3.62);

  \node[text=source,align=center,font=\scriptsize] at (5.86,.75)
    {统一与守边依赖区域资源进入中枢；实际征收与运输并非单向无阻};
\end{tikzpicture}

\smallskip
\small\textit{图\quad 北宋初年区域资源、运输与国家能力关系示意：据官修纪传《宋史》卷1、卷2〈太祖本纪〉及后蜀、南汉、南唐世家，南宋编年汇集《续资治通鉴长编》乾德、开宝诸年条，后世辑存的制度汇编《宋会要辑稿·方域》，并参照开封北宋东京城遗址的水系考古资料，说明东京、边防、四川、江南和岭南之间的资源接收与调度关系。它不表示税额、运量、行政执行速度或某一地区必然服从中枢；图中区域边界、水道和资源、军需、信息路线均为约略示意。}
\end{center}


976年太祖去世，赵光义即位为太宗。即位日期可由《宋史》两朝本纪与《长编》开宝九年条确认；关于继承的“烛影斧声”却是后起传闻，现存材料并未使皇子与兄弟之间的安排完全透明。把这一不透明处改写成确定的宫廷阴谋，只会用后世兴趣覆盖当时的政治难题。太宗所继承的，是已扩大的财政与军队，也是尚未结束的河东割据和北方边患。

\section{河东之后，没有缓冲地带}

979年五月，宋灭北汉，河东归入宋朝，通常被视为五代十国割据大势的收束。《宋史》太宗本纪、北汉世家和《长编》太平兴国四年条足以确证灭国这一结果；北汉长期倚重辽的援助这一事实，也说明太原并非孤立城池。宋军取得河东后，宋与辽之间原先的缓冲政权随之消失。统一把军队和赋税推到了更远的北方，也把边界变成直接相接的战场与道路。

太宗旋即北攻幽州，在高梁河失利南撤。宋方的《宋史》《长编》和辽方的《辽史》景宗本纪都记下这场挫败，却在兵力、皇帝是否受伤、追击路线及功过归属上并不相同。因而高梁河不能只写成一方“失机”或另一方“天佑”的故事：宋军希望趁河东既下夺取燕云，辽军则能依托近畿、骑兵机动与既有防御应对。战败使速取燕云的设想落空，也迫使东京重新计算北部防线需要多少兵员、粮饷和指挥协调。

986年的雍熙北伐又以三路军北进，终至受挫撤军。三路并进的构想反映出宋廷仍想以协同动作牵制辽军；宋、辽两部史书与《长编》雍熙三年条同时提醒我们，诸路军的战果、军数和将领责任都经过战后叙事塑形。可以把握的后果是，北伐未能改变燕云的格局，守势边防和军费压力反而更深地进入中枢决策。此前限制将领独立兵权的制度，避免了五代式军头坐大，却也使远方战役更依赖来自东京的调度、文书和供给；这不是单一制度造成的胜败，却构成宋初军事协调所承受的约束。

997年赵恒继位为真宗，承接的正是这样的北边与财政。1004年冬，辽圣宗和萧太后南下，真宗至澶州前线，双方在军事对峙中转入议和。宋方的真宗本纪和《长编》景德元年冬条，与《辽史》圣宗本纪共同构成这一事件的基本证据；至于谁在战场上占有绝对主动、谁以何种言辞劝动皇帝，则多被后来的史家赋予戏剧性的形状。澶州的选择并不是抽象的怯战或勇战，而是两国都在机动、城防、补给和内部政治压力下衡量继续作战的成本。

1005年盟约落实，宋输送银绢，双方以兄弟之国交往。宋方记载中的“岁币”语言常被后世用来衡量屈辱或胜负；《宋会要辑稿·蕃夷》所存盟约条和《辽史》的并读，更适合把它放回两国关系的制度语境：银绢、使节、边市与防务被编入一套长期的停战安排，不能单凭一方的朝贡称谓推出单向统属。1008年的封禅和相关大礼，显示宋廷又尝试把边境安定转译为皇权与天命的语言；《宋史》真宗本纪、《长编》和《宋会要辑稿·礼》可见礼仪的举行与制度记录，却不能把祥瑞叙事当作事实的验证。

\phantomsection\label{fig:vol04:northern-song:song-liao-chanyuan-frontier}
\begin{center}
\begin{tikzpicture}[x=.88cm,y=.88cm,font=\small]
  \fill[paper] (0,0) rectangle (11.7,6.35);

  \fill[source!20] (.5,3.7) -- (11.2,3.45) -- (11.25,5.9) -- (.7,5.85) -- cycle;
  \node[align=center] at (9.5,4.95) {辽\\燕云、骑兵机动与北方近畿};
  \fill[dynasty!22] (.55,.55) -- (11.25,.7) -- (11.2,3.25) -- (.55,3.45) -- cycle;
  \node[align=center] at (2.0,1.72) {宋\\河北防务与东京中枢};

  \draw[very thick,color=timeline!85] (.7,3.52) .. controls (3.05,3.88) and (5.7,3.15) .. (8.0,3.55)
    .. controls (9.25,3.75) and (10.2,3.4) .. (11.15,3.55);
  \node[draw=timeline!80,fill=paper,rounded corners=2pt,align=center,inner sep=3pt] (line)
    at (5.95,3.6) {盟约维系的\\边境接触地带};

  \fill[timeline] (6.0,2.55) circle (.13);
  \node[below,align=center] at (6.0,2.45) {澶州\\1004年对峙};
  \fill[source] (3.4,4.6) circle (.13);
  \node[above,align=center] at (3.4,4.68) {燕云方向};
  \fill[dynasty] (9.55,1.65) circle (.13);
  \node[below,align=center] at (9.55,1.55) {东京\\开封};

  \draw[->,ultra thick,color=source!85] (3.65,4.45) .. controls (4.55,4.0) and (5.05,3.25) .. (5.75,2.68);
  \node[text=source,align=center,font=\scriptsize] at (4.2,3.4) {1004年辽军\\南下方向};
  \draw[->,ultra thick,color=dynasty] (9.3,1.72) .. controls (8.0,1.9) and (7.05,2.2) .. (6.25,2.52);
  \node[text=dynasty,align=center,font=\scriptsize] at (7.9,2.4) {宋廷赴澶州\\及河北动员};

  \draw[<->,thick,dashed,color=timeline!90] (5.45,3.06) -- (6.6,3.08);
  \node[text=timeline!90!black,above,align=center,font=\scriptsize] at (6.02,3.08)
    {使节、盟约与\\银绢输送关系};
  \node[draw=source!70,fill=paper,rounded corners=2pt,align=center,inner sep=3pt] at (9.3,4.8)
    {和平依赖\\防务、谈判、边市\\与地方管控并行};
  \node[text=source,align=center,font=\scriptsize] at (2.0,.85)
    {1005年后并非无边防的“固定国界”};
\end{tikzpicture}

\smallskip
\small\textit{图\quad 澶渊盟约前后的宋辽边境关系示意：据宋、辽官修纪传《宋史》卷5〈真宗本纪一〉与《辽史》〈圣宗本纪〉、南宋编年汇集《续资治通鉴长编》景德元年冬条，以及后世辑存的制度汇编《宋会要辑稿·蕃夷》辽国盟约条，呈现1004年澶州对峙、1005年盟约与防务、使节、银绢和边市之间的关联。双方战场主动性、银绢数额及各方称谓须回到文本语境辨析。图中政权边界、盟约接触地带与军队、使节和输送路线均为约略示意。}
\end{center}


盟约也没有把边界变成无人之地。1010年辽军再入宋境，宋廷加强河北边防，同时维持盟约渠道。《宋史》《辽史》与《长编》对破约和战果各有立场；能够看见的是，和平依赖的从来不只是纸上的誓词，还包括堡寨、侦察、地方管控、军队供给与使节往返。北宋最初数十年的国家能力，正在这种并行中成形：开封把资源向中枢汇集，统一战争把新的区域纳入账目与道路，北方压力又迫使朝廷不断校验它能调动多少人、粮和信息。它解决了五代以来最迫切的分裂与政变风险，却也把一套昂贵而持续的边防财政，留给了此后更长的宋辽关系。

\section{摄政之后，西北成为一张持续的账单}

1022年真宗去世，年幼的赵祯即位，刘太后临朝。继位和摄政可由《宋史》仁宗本纪、章献明肃刘皇后传与《续资治通鉴长编》乾兴元年条互证；不过，把近十年政务逐项归为太后个人意志，正是后世宫廷叙事容易造成的错觉。她所维持的是一个须同时支付河北防务、辽宋盟约和官僚日常运转的中枢。1033年刘太后去世，仁宗亲政，改变的是决断与人事的重心，并不是一夜之间出现了不受辅臣、台谏和内廷约束的皇帝。

西北的变化使这种继承很快受到检验。1034年李元昊改国号、年号，1038年称帝，意味着党项首领不再满足于宋朝册封秩序。《宋史》夏国传、《长编》与《宋会要辑稿·蕃夷》能确定宋廷所面对的外交转折；它们却主要记录宋朝如何命名对手。西夏的建国礼、内部接受和财政基础，不能只从敌国战报反推，仍须由黑水城文书、碑刻和区域考古补足。对开封而言，关键并非一句称号是否“僭”，而是河套、河西道路、堡寨与水源附近，出现了一个能独立调动资源的政权。

1040年的三川口、1041年的好水川和1042年的定川寨，连续打破宋军速胜的期待。宋方本纪、范雍传、任福传和《长编》保留了败讯，却把将领责任、诱敌经过与伤亡数字包裹在战后追责中；不宜把奏报的军数当作已被核实的战场规模。三场战争共同呈现的约束较为清楚：黄土高原上的行军、给水、粮道和情报，限制了从东京发出的命令；西夏军队则能够把机动和地形结合为防御。于是边防不只是地图上一条线，而是一串需要驻兵、修寨、运粮和轮换的支点。

\section{改革的尺度与边境的尺度}

这种压力解释了1043年庆历新政何以在西夏战争中提出。范仲淹、富弼等人针对选任、考课与学校提出措施，希望重整官僚的入口和责任；《宋会要辑稿·职官》庆历条与《长编》可以追索政策的提出。它们证明诏令与议论，却不足以证明州县已经同样执行。到1045年，主要倡议者相继离开中枢，改革的整套推进停止。若把这段经历简化为某一派“胜”或“败”，便看不见考课、荐举和既得利益如何在不同层次制造阻力；它留下的是一批未消失的问题和一份日后反复被援引的改革记忆。

1044年宋夏和议则表明，战争与改革都不能脱离财政的限度。双方以赐予、互市和停战安排降低边境消耗。《宋史》仁宗本纪、夏国传和会要的和议条足以确认框架；“赐予”这一宋方用语不能自动换算为单向臣属，条款、称谓和实际执行应分开处理。和平没有取消堡寨、侦察和贸易摩擦，却让朝廷可以重新配置西北支出。这是以军备支撑的停战，而不是没有代价的安宁。

边疆也不只在西北。1052年侬智高在岭南与交趾边境起兵，活动及邕州、广州附近；次年狄青率军平定。《宋史》侬智高传和狄青传提供宋方的行动线索，但其族群标签和“平定”后的范围，都不等于当地社会的自我分类和全部恢复。广源一带的首领原在宋、交趾与山地网络之间争取资源与承认，朝廷调兵虽重获主要城镇，却没有消除跨境道路、地方武装和交换关系。仁宗朝的边防由此显出双重面貌：中枢可以集中军费和将领，却不能把边地的多重联系压缩成一纸州县名册。

\section{从继统之争到富国强兵}

1063年仁宗去世，英宗继位。仁宗无存活皇子的继承安排，需要宗室、后族、宰执和礼官共同承认；《宋史》两朝本纪和《长编》嘉祐八年条可确认这一转换。1065年的濮议随即把英宗生父濮王应获何种追尊推到朝廷中心。它不是可以被还原成纯粹礼学或纯粹党争的一次争吵：宗法名分怎样写入国家礼制，会决定谁有资格解释皇统，也会重新结连台谏和宰执。

1067年神宗即位，财政、兵役和边防的未决问题进入新阶段。次年王安石入朝，神宗与他讨论变法方向；《宋史》神宗本纪、王安石传和《长编》提供了人事与论议的骨架。把其后的全部政策都写成一人设计、全国照办的蓝图，会遮蔽皇帝、制置机构、地方官和胥吏各自的作用。新法首先是一套试图把财政来源、基层劳役与军事压力重新接合的制度实验。

1069年青苗法颁行，官府在青黄不接时以贷款介入借贷；1070年前后免役法在部分地区推进，以雇募和免役钱改造差役；1071年保甲法开始以户籍编组治安、训练与互保。《宋史》食货志、兵志、《宋会要辑稿》与《长编》能分别确认这些制度文本和推进节点。它们不能证明同一利率、同一征收方法或同一训练密度覆盖全国。对小民而言，政策的实际面貌取决于户等认定、地方官考成、书手催科和市场中的货币；正是在这些环节，救济、加派、逃避与争讼可能并存。

\section{新法走向边境，也走进文字}

1072年市易法设市易务，以官本介入城市商货和信贷。它说明改革并未只瞄准乡村：开封等城市已是大宗货物、借贷与官府需求交会之处。会要和食货志所记的机构、资金并不能单独证明物价稳定或商人盈亏；还须由区域性市场材料检验。制度争论的核心也不只是“官”与“商”的抽象对立，而在于谁掌握货物信息、何时催收本息、谁承担价格波动。

1074年旱灾、政策批评和人事冲突交织，王安石首次罢相。灾异在政治语言中常被用作责难政府的证据，不能据此断言旱灾直接造成政策逆转。此后改革并未终止，而是更依赖神宗与不同官僚网络的支持。1079年苏轼因诗文被控讥讽朝政，乌台诗案显示台谏、法司和文人如何以诗句、奏议和“结党”语言划定忠诚边界；《宋史》苏轼传与《长编》可证案件与处分，却不容把每句诗的政治含义写成唯一解释。

\subsection{交读窗口：黄州的水面与乌台案的文书}

元丰五年（一〇八二）秋，仍在黄州的苏轼写《前赤壁赋》。这是以舟游、问答和主客转换铺写的散赋，不是申诉状或判牍；“寄蜉蝣于天地，渺沧海之一粟”把遭贬者置于江水与时间的尺度中。它之所以能照亮乌台诗案之后的政治文化，不在于替案件供给一条新事实，而在于让人看见一位被处分的官员如何把失意、友游与经典修辞转换为可公开流传的文字。

案件的日期、牵连的奏议与处分，仍应以《续资治通鉴长编》元丰二年条、《宋史》〈苏轼传〉及可追查的《宋会要辑稿·刑法》条目为脊梁；赋中“客”与“苏子”的声音也不能逐一还原成法司、台谏或黄州居民的意见。钱穆所关心的制度与士人之间如何彼此制约，在此可转化为一个具体问题：案件后的言说空间由谁界定？吕思勉式的社会史追问则要求把诗文名士与催科、保甲、城寨中的人分开观察。地方墓志、题名和文集的保存尤偏向识字精英，故《前赤壁赋》可以说明一段文学姿态，不能量度新法执行，亦不能代表所有被政令卷入者的感受。

边疆的账目同样没有消失。1075年李朝军队攻入邕、钦、廉州，1076--1077年郭逵率军南征至富良江一线后撤。宋、越材料都承认战争，却在战果、疫病与损失上保留各自叙事位置；热带道路和补给的限制，比单一的勇怯判断更能解释未能长期占领。1080年永乐城失守，1081年五路伐夏未获预期战果，城寨的选址、供水、运输和多路协同再度暴露为军事问题。1082年元丰改制重整官制名目与机构关系，意在处理权责交错；复古的官名不等于实际职能立刻稳定。新法的代价与边防的代价，都仍由中枢的文书、财政和地方执行来承担。

\section{制度可以废改，政治记忆并不随之消失}

1085年神宗去世，哲宗即位，高太后临朝。年幼皇帝与高度争议的新法体系，使摄政并非简单的“复旧”，而是必须逐项判断哪些法令停止、替换或保留。1086年元祐政府废罢、调整青苗、免役等多项新法。《宋会要辑稿·食货》的废改条与《长编》可追踪中枢决定，但不同制度的废止时间、地方保留与替代措施并不同步。把一纸诏令当作全国同时翻转，恰好会重复新旧两派各自最方便的叙述。

1093年高太后去世，哲宗亲政；1094年绍圣政府复行多项熙宁新法，并贬逐元祐官员。继承、摄政终结和人事转向的事实，可由《宋史》本纪、传记和会要相互核对；将其全归于皇帝个人意志，则忽略宰执、言官和外放官员形成的政策网络。制度在这里也成为官员履历的一部分：拥护或反对某项法，逐渐会被记录为可供任用、贬谪和结交的身份。

1097年继续编列、处分元祐异议官员，尤其显示这种转换如何把政策分歧变成行政分类。党人名单的版本、人数和名单中人的实际主张并不完全一致；《宋史》所用的道德化传名更不能直接等同于每位官员的政治生涯。可确定的是，名册、贬所和人事禁限使异议具有了制度化的地理与记忆。北宋中枢并非从此不能运转，却愈来愈难在政策分歧之外保留稳定的合作空间。

\section{宫廷的摇摆与北方的新行动者}

1100年哲宗无嗣而死，徽宗即位。端王被选择的具体过程材料有限，后出的预言与人物评语不能替代可核验的宫廷程序。1102年蔡京再入中枢，1104年元祐党人碑公开标列异议者；1106年朝廷又一度调整崇宁政策、宽释党人，1107年蔡京复相。四个节点见于《宋史》徽宗本纪、蔡京传和会要相关条目，足以说明人事与政策的反复，却不支持把所有晚期政务或全部危机归咎于一人。党籍、学校、财政和任官被连在一起，令朝廷的短期调整更难摆脱既有依附网络。

1115年完颜阿骨打建金、反辽，北方格局因此重组。《金史》太祖本纪、辽天祚帝本纪和宋徽宗本纪可以交叉确证建国与战争的基本序列；宋廷初获何种情报、何时形成判断，仍须细读使节材料。女真集团的兴起既来自辽统治下的压力，也依赖完颜联盟的组织能力，不宜写成宋朝外交所能任意操纵的机会。

1118年宋廷经海路和边境渠道接触金方，1120年达成海上之盟，共攻辽并商议燕云。宋、金交聘表与《三朝北盟会编》都保存谈判和盟约线索，但条款传本、岁币转换和燕云范围有异文。此盟改变了宋辽之间已延续多年的停战框架：宋希望借金军收复燕云，金则需要南方牵制辽军。它不是一份自动保证收复的契据，而是把宋的军费、作战能力和战后边界安排置于更严厉的检验之下。

\section{燕京的交接不是收复的终点}

1121年宋军在两浙平定方腊后，1122年童贯等依盟约攻辽失利；同年金军攻取辽南京燕京，辽朝北撤。宋、辽、金三方记载对军纪、抵抗和失利缘由各有立场，城内破坏与居民处境也不能由任一方的胜利叙述概括。可以把握的是，宋未能独力取得燕京，金因而掌握交接条件。

1123年金将燕京及部分地区交给宋，宋承担岁币与边防成本。名义上的移交不等于实际控制：范围、户口、城防和驻军须逐项核对。新前沿既脆弱又昂贵，宋金互信迅速下降。1123--1124年平州守将张觉降宋，宋廷接纳后成为争端；决策过程和张觉的地方基础并不完整，但双方围绕降将与缓冲地带的争执确为1125年金军南侵的直接借口之一。联盟曾承诺的燕云，在此已由目标变成一条难以供给、难以解释、也难以防守的边界。

\section{东京被围：决策的时间比城墙更短}

1125年冬金军分路南侵，宋金战争全面爆发。张觉事件、燕云交接的矛盾和金方对宋军能力的判断共同构成背景，但不应由一个借口解释全部战争。《宋史》徽宗、钦宗本纪，《金史》太宗本纪和《三朝北盟会编》构成基本证据链；军数、城池控制和宋廷何时得知敌情，仍须逐城、逐月核验。

1126年正月，金军第一次围汴京。城内必须同时处理守城、勤王与和议，三者都需要时间、粮饷和可靠消息，而这些条件并不总能同时具备。同月徽宗传位钦宗，自称太上皇。传位的事实清楚，却不能简单写成个人逃避：围城压力和宫廷程序共同迫使皇室寻找新的谈判与统治重心。其结果反而是双重权威，钦宗须在太上皇、宰执与各路勤王军之间作出选择。

第一次解围没有改变金军的战略优势。赔款、割地与人事争论使临时的喘息更难转为一致防务。

\begin{event}{\eventanchor{NS-1127-JINGKANG-FALL}{靖康汴京陷落}（一一二七年四月；金与宋；汴京；城陷可核，俘掠数额不能精算）}
《宋史·钦宗本纪》《金史·太宗本纪》与《三朝北盟会编》可互校城陷及徽、钦二帝被掳的核心事实。俘虏、财物和暴行的具体数字带有战时与后出叙事的放大，不能相加为确定统计，也不能代替城内外不同人群的经历。
\end{event}

1127年四月汴京陷落，徽、钦二帝及大量宗室、官民被掳，北宋中央政权终结。宋金两方史书和北盟材料都可确认城陷与俘虏的核心事实；人数、财物和暴行的具体数额，多带有战时与后出叙述的放大，不能相加为确定统计。

靖康不是一座城市突然消失的瞬间。它意味着仓储、官署、家族、俘虏迁徙和道路网络被战争重新切断或接续；幸存宗室与官民南渡，南宋建国和宋金长期战争由此开始。北宋留下的中央集权、财政技术与市场网络没有全数消散，但它们再也不能由开封这一个中枢统摄。下一阶段的历史，须在南方行在、北方金朝与仍在流动的人群之间重写。

